\section{总结与延伸}

\subsection{本讲核心要点}

\begin{enumerate}
    \item \textbf{SDE 统一框架}：离散的 DDPM 和 SMLD 都可以自然地推广到连续时间 SDE。前向过程由漂移系数 $f$ 和扩散系数 $g$ 完全确定。

    \item \textbf{VE-SDE 与 VP-SDE}：
    \begin{itemize}
        \item VE-SDE（对应 SMLD）：$f = 0$，$g = \sqrt{\mathrm{d}[\sigma^2]/\mathrm{d}t}$，方差单调递增；
        \item VP-SDE（对应 DDPM）：$f = -\frac{1}{2}\beta(t)\mathbf{x}$，$g = \sqrt{\beta(t)}$，方差保持恒定。
    \end{itemize}

    \item \textbf{反向 SDE}：由 Anderson 定理直接给出，只需额外知道 score function。训练一个 score 网络后，前向和反向过程就完全确定了。

    \item \textbf{Variational Diffusion Models}：通过 $(\alpha_t, \sigma_t)$ 参数化，在 SNR 单调递减的条件下，定义了一个无限的扩散模型家族。VE-SDE 和 VP-SDE 只是其中两个实例。

    \item \textbf{Probability Flow ODE}：每个 SDE 都有一个对应的 ODE，产生相同的边际分布但具有确定性轨迹。这为高效采样和精确似然计算提供了工具。
\end{enumerate}

\begin{importantbox}{连续时间视角的价值}
连续时间 SDE 框架的价值不仅在于"统一"已有方法。更重要的是，它为我们提供了设计新型扩散模型的系统性方法论——只需定义 $f$ 和 $g$（或等价地，$\alpha_t$ 和 $\sigma_t$），剩下的（反向过程、ODE、训练目标）都自动确定。这种"plug and play"的灵活性催生了大量后续工作。
\end{importantbox}

\subsection{后续方向}

本讲建立的连续时间框架是后续多个重要主题的理论基础：

\begin{itemize}
    \item \textbf{高效采样器}：基于 Probability Flow ODE 的高阶数值求解器（DPM-Solver, DEIS 等），可以在 10--15 步内生成高质量样本；
    \item \textbf{Flow Matching / Rectified Flow}：一种更简洁的连续时间生成模型训练方法，可以看作对 Probability Flow ODE 的直接参数化；
    \item \textbf{Consistency Models}：通过强制 ODE 轨迹上不同时刻的映射一致，实现单步生成；
    \item \textbf{精确似然评估}：利用 ODE 和 Hutchinson trace estimator 计算对数似然。
\end{itemize}

\subsection{拓展阅读}

\begin{itemize}
    \item Song, Y., et al. (2021). ``Score-Based Generative Modeling through Stochastic Differential Equations.'' \textit{ICLR 2021.} ——本讲的核心参考文献。
    \item Anderson, B. D. O. (1982). ``Reverse-time diffusion equation models.'' \textit{Stochastic Processes and their Applications.} ——反向 SDE 的经典理论。
    \item Kingma, D. P., et al. (2021). ``Variational Diffusion Models.'' \textit{NeurIPS 2021.} ——Variational Diffusion Models 框架的提出。
    \item Karras, T., et al. (2022). ``Elucidating the Design Space of Diffusion-Based Generative Models.'' \textit{NeurIPS 2022.} ——系统性探索连续时间扩散模型设计空间的重要工作。
\end{itemize}

%% --- Fin del contenido --- %%

\end{document}
